\begin{afterword}
\summaryhead

The post continues the previous day's claim that rationality is systematized winning. To the objection that reasonable people do not always win, it answers first that lottery winners are a product of selective reporting. To the stronger objection, that overconfident founders and religious believers seem to do better, it concedes a point: the optimality theorems hold only for perfect reasoners, and a step toward rationality need not bring more winning. There can be a valley.

The author gives three personal reasons to take the step anyway: the author's work punishes small errors, refusing one step forfeits all the later ones, and a deception once seen cannot be unseen. The author then reports that experience suggests that, once a person stops making other mistakes, more rationality does make them better off. This is plain in answering confusing questions, but not yet shown for motivation or happiness. Two reader comments about losing a girlfriend after deconverting illustrate the valley. The post argues that arguments against losing illusions usually consider one step only, and ends with a joke telling readers who want their illusions to stop reading.

\responsehead

The post's concession is its strongest part, and it is correct. There is no theorem that each step toward ideal reasoning improves results, and the post admits it plainly. Economists had proved a close relative of the point in the theory of the second best (1956): when one condition of the optimum cannot be met, meeting the others may make things worse. The book itself warned in 2007, in ``Knowing About Biases Can Hurt People,'' that more knowledge can hurt someone who is irrational to start with.

What the post sets against the valley is weaker, and the post mostly knows it. Its three reasons are the author's own, offered after ``I can't necessarily answer for everyone.'' Its general answer is based on the author's experience, hedged (``does seem to indicate'') and not described. In the field the author has ``most intensely practiced,'' the claim that an optimist would be ``stomped into the ground'' by ``a casual student'' comes with ``seems'' and no example.

On motivation and happiness the post is candid. It cannot guarantee that losing illusions helps, and it states a hope and a belief, marked as such. But the question was not new. Psychologists had argued it for two decades: Taylor and Brown (1988) held that positive illusions foster mental health, and Colvin and Block (1994) replied that this remained unproven. The post refers to ``the surveys'' without citing one. Two later findings bear on the picture. A 2019 meta-analysis found self-enhancement linked to better personal adjustment in ordinary people, which fits the objection the post answers, though it says nothing about the best a person could achieve. Some studies that separate weak from strong belief suggest worse mental health in the middle than at either end, which fits the post's valley.

The author's remark that no systematic art of happiness is known to the author should be read beside the author's own statement, eight days earlier, that cognitive-behavioral therapy had been ``experimentally shown to yield real improvements.'' The claim is hedged with ``If,'' and CBT corrects distressing beliefs more than comforting ones, so this is information for the reader rather than a contradiction.

The closing conditions include being signed up for cryonics, treated as the rational choice. The argument for that is in a post the book does not include.

In short: a correct and candid concession that partial rationality can leave a person worse off, answered with the author's reasons, experience and hopes, and with no study on either side.
\end{afterword}
